% UTM individual report (laboratory work, course or semester report),
% with the same layout as the mandatory ELSE Word template.
%
%   1. Fill in the cover fields below.
%   2. Replace the red example text; everything is in this one file.
%   3. Put images in figures/ and include them by name: \includegraphics{my-diagram}.
%   4. Compile from this folder:  tectonic main.tex   (produces main.pdf)
%
% Red text in the PDF comes from \placeholder{...}: replace all of it before submitting.
% Do not edit utmreport.cls; it holds the margins, fonts and heading styles.
\documentclass{utmreport}
\selectlanguage{english}

% ---- Cover -----------------------------------------------------------------
\institution{Technical University of Moldova}
\faculty{Faculty of Computers, Informatics, and Microelectronics}
\department{Department of Software Engineering}
\studyprogramme{Software Engineering Study Programme}
\documenttype{Laboratory work No. 1}
\reporttitle{\placeholder{Title of the laboratory work}}
\student{\placeholder{Name SURNAME}}{\placeholder{FAF-24x}}
\supervisor{\placeholder{Name SURNAME}}{\placeholder{univ. assist.}} % {name}{title}; title may be empty
\cityyear{Chișinău}{2026}

\begin{document}
\makeutmreporttitlepage
\tableofcontents

% ---- Abstract: one page, 5-7 lowercase keywords ----------------------------
\begin{utmabstract}[Abstract]
\placeholder{Summarize the task, the method, the result and the main conclusion in one paragraph.}

\utmkeywords{\placeholder{requirements, methodology, implementation, evaluation, documentation}}
\end{utmabstract}

% Optional: required when the report uses abbreviations or English terms.
% \begin{utmabbreviations}
% \item[API] Application Programming Interface
% \end{utmabbreviations}

% ---- Introduction: purpose of the work and the task conditions -------------
\utmintroduction
\placeholder{State the purpose of the laboratory work and the task conditions.}

\placeholder{The objectives of the work are:}
\begin{utmenumerate}
\item \placeholder{explaining the problem and the method;}
\item \placeholder{implementing and testing the solution;}
\item \placeholder{drawing conclusions supported by the results;}
\end{utmenumerate}

% ---- Chapters: each starts on a new page with at least one paragraph ---------
\chapter{Methodology}
\placeholder{Explain how the requirements, the implementation and the checks are connected before the first subchapter.}

\section{Model and notation}
\placeholder{Equation \ref{eq:sum} defines the sum of $n$ values $x_i$; $S$ denotes their total:}
\begin{equation}\label{eq:sum}
S=\sum_{i=1}^{n}x_i.
\end{equation}

\section{Implementation}
% Courier New 10 pt, single spacing. Code longer than half a page goes in an appendix.
\placeholder{Short code belongs beside its explanation; the complete program is in Appendix \ref{app:code}. The following loop calculates the same sum:}
\begin{lstlisting}
int sum = 0;
for (int index = 0; index < values.Length; index++)
{
    sum += values[index];
}
\end{lstlisting}

\chapter{Results}
% Never start a page with a table or figure: introduce it in a short paragraph first.
\placeholder{Table \ref{tab:example} compares the values before and after the change; the values here are illustrative, not measurements.}
\begin{table}[htbp]
\caption{\placeholder{Comparison of the results}}\label{tab:example}
\centering
\begin{tabularx}{0.9\textwidth}{|Y|c|c|}\hline
\textbf{Indicator} & \textbf{Initial value} & \textbf{Final value}\\\hline
Processing time, min & 42.50 & 31.25\\\hline
Errors, count & 7 & 2\\\hline
\end{tabularx}
\end{table}

\placeholder{According to Figure \ref{fig:example}, the input is processed in one step. Cite the source of a reused or adapted figure \cite{else-model}.}
\begin{figure}[htbp]
\centering
\includegraphics[width=0.45\textwidth]{example-diagram}
\caption{\placeholder{Processing flow}}\label{fig:example}
\end{figure}

% ---- Conclusions -----------------------------------------------------------
\utmconclusions
\placeholder{Relate the results to the objectives, explain the limits and identify the next steps.}

% ---- References: numbered in order of first citation, each cited with \cite --
%   article  AUTHOR 1, A.B.; AUTHOR 2, C.D. Title. Abbreviated Journal Name Year, Volume, pages, doi.
%   book     AUTHOR 1, A. Book Title, 3rd ed.; Publisher: Location, Country, Year; pp. 154-196.
%   website  Title of Site. Available online: URL (accessed on Day Month Year).
\begin{thebibliography}{9}
\bibitem{else-model} Technical University of Moldova. \textit{Mandatory report template} (ELSE model). Available online: \placeholder{URL} (accessed on \placeholder{Day Month Year}).
\end{thebibliography}

% ---- Appendices: "APPENDIX A – Title", each on a new page --------------------
\appendix
\chapter{Source code}\label{app:code}
\placeholder{Complete program listing referenced in the main text.}
\end{document}
